% Prijava teme seminarske naloge -- NPMP 2026/2027, kvantno procesiranje
% Skupina 20. Prevajanje: pdflatex predlog_teme.tex
\documentclass[10pt,a4paper]{article}
\usepackage[utf8]{inputenc}
\usepackage[T1]{fontenc}
\usepackage[slovene]{babel}
\usepackage{mathptmx}
\usepackage{courier}
\usepackage[a4paper,margin=1.7cm,top=1.4cm,bottom=1.4cm]{geometry}
\usepackage{amsmath}
\usepackage{booktabs}
\usepackage{enumitem}
\usepackage[hidelinks]{hyperref}
\setlist{nosep,leftmargin=1.2em}
\setlength{\parindent}{0pt}
\setlength{\parskip}{4pt}
\pagestyle{empty}
\sloppy

\begin{document}

{\large\textbf{Prijava teme seminarske naloge -- NPMP 2026/27, kvantno procesiranje}}\\
\textbf{Skupina 20:} Dominik Krašovec, Erik Radovičevič \hfill Ljubljana, 2.\,11.\,2026

\medskip
\textbf{Naslov:} \emph{Stanja GHZ na procesorju IBM}

\textbf{Raziskovalno vprašanje.} Do koliko kubitov $N$ lahko na šumnem procesorju IBM Heron (ibm\_fez, 156 kubitov)
pripravimo stanje GHZ, da meritev še pokaže prepletenost? Kako izbira kubitov po povezavah čipa (veriga proti drevesu)
vpliva na to mejo?

\textbf{Stanje GHZ.} $N$ kubitov v stanju
$\lvert\mathrm{GHZ}_N\rangle=\tfrac{1}{\sqrt2}(\lvert0\rangle^{\otimes N}+\lvert1\rangle^{\otimes N})$ je prepletenih:
stanja ni mogoče zapisati kot stanja vsakega kubita posebej. Stanje pripravimo z vrati H na prvem kubitu in $N-1$ vrati
CNOT po povezavah čipa. V idealnem primeru so po vratih H vsi nizi izidov s sodo pariteto enako verjetni.

\textbf{Hipoteza.} Na idealnem simulatorju je priprava točna. Šum (napake dvokubitnih vrat, čakanje kubitov, napake branja)
kakovost stanja z večanjem $N$ zmanjšuje. Na drevesu povezanih kubitov je vezje plitvejše kot na verigi, zato pričakujemo,
da drevo pri istem $N$ obdrži več prepletenosti. Stanja GHZ so na procesorjih IBM že pripravili s 27 kubiti; zanima nas,
koliko pridemo na novejšem procesorju ibm\_fez z vezjem, ki ga lahko sestavi začetnik, in z brezplačnim procesorskim časom.

\textbf{Metrike.}
\begin{itemize}
\item Lastnosti vezja pred prevajanjem in po njem: število kubitov, globina, število dvokubitnih vrat, trajanje.
\item Primerjava idealnega simulatorja, simulatorja s šumnim modelom procesorja in procesorja ibm\_fez pri istem $N$.
\item Izbira kubitov po objavljenih napakah procesorja (veriga in drevo).
\end{itemize}

\textbf{Načrt poskusov.}
\begin{enumerate}
\item \emph{Vezje:} priprava stanja GHZ z vrati H in $N-1$ vrati CNOT po povezavah čipa (brez dodanih vrat SWAP).
\item \emph{Idealni simulator} (Qiskit Aer) in \emph{simulator s šumnim modelom} (FakeFez).
\item \emph{Procesor ibm\_fez:} naloge z vnaprej določeno zgornjo mejo trajanja (paket Open Plan: 600~s v 28 dneh).
\item Primerjava veriga / drevo in razprava o omejitvah (izbrani kubiti, meritev enega dne).
\end{enumerate}

\textbf{Procesorski čas.} Paket Open Plan omogoča 600~s v 28 dneh; vsaka naloga ima vnaprej določeno zgornjo mejo
trajanja. Pred oddajo bomo porabo ocenili in je ne bomo prekoračili.

\textbf{Orodja.} Python~3, Qiskit~2.6 (vezja, prevajanje), Qiskit Aer 0.18 (idealni in šumni simulator), qiskit-ibm-runtime
0.51 (SamplerV2), NumPy, Matplotlib, pytest; IBM Quantum Open Plan; \LaTeX{} (predloga zbornika); GitHub.\par

\textbf{Razdelitev dela.} Oba člana (Dominik in Erik) delava na vseh delih naloge: gradnja vezij in izbira kubitov,
simulacije, zagon na IBM in beleženje ID naloge ter porabe časa, analiza, pregled literature, poročilo, prosojnice in
repozitorij (README). Vsak del napiše ali izvede eden, drugi ga pregleda in neodvisno ponovi.

\end{document}
